\begin{center}\textls[100]{b. Die späteren Beweise des Fundamentaltheorems (1808—1818).}\end{center}
\begin{center}\bfseries 18.\end{center}
\noindent\hspace*{\parindent}Nach den D. A. veröffentlichte \textsc{Gauss} zunächst die drei Abhandlungen: \textit{Theorematis arithmetici demonstratio nova} (Comm. Gotting. \textbf{16}, 1808, vorgelegt 15. Januar); \textit{Summatio quarundam serierum singularium} (Comm. Gotting. rec. \textbf{1}, 1811, vorgelegt am 24. Aug. 1808); \textit{Theorematis fundamentalis in doctrina de residuis quadraticis demonstrationes et ampliationes novae} (ebendas. \textbf{4}, 1818, vorgelegt am 10. Febr. 1817). Sie enthalten vier neue Beweise des Fundamentaltheorems. Die zweite von ihnen ist, wie die \textit{Disquisitionum circa aequationes puras ulterior evolutio}, deren Fortsetzung sie ursprünglich zu bilden bestimmt war (s. W. II, S. 265, Bemerkungen), unmittelbar aus der Kreisteilung hervorgegangen. Dort finden sich in art. 356 der D. A. bereits die sogenannten \textsc{Gauss}schen Summen, denen man zusammenfassend die Gestalt
\[
W_q=\sum_{i=0}^{n-1}r^{qi^2}
\]
geben kann, wo
\[
r=\cos\frac{2\pi}{n}+\sqrt{-1}\sin\frac{2\pi}{n}
\]
gesetzt ist, und deren Wert für den Fall einer ungeraden Primzahl \(n\), bis auf das Vorzeichen einer Quadratwurzel genau, dort angegeben ist. Auch
\newpage\SourceLeaf{53}\ArithmeticHead{46}{L}{BACHMANN, ÜBER GAUSS’ ZAHLENTHEORETISCHE ARBEITEN.}
\noindent ohne dieses letztere zu bestimmen, erkannte \textsc{Gauss} hier schon Mitte Mai 1801 (Nr. 118 T.) einen neuen Weg zum Beweise des Reziprozitätsgesetzes. Aber die Bestimmung des Vorzeichens selbst gestaltete sich zu einer ebenso reizvollen wie schwierigen Aufgabe, die zu bewältigen vier Jahre erforderlich waren. An mehreren Stellen (W. II, S. 16, 156; Nr. 123 T.), mit besonderer Lebhaftigkeit aber in einem Briefe an \textsc{Olbers} vom September 1805 (W. X 1, S. 25) hat \textsc{Gauss} die Mühen geschildert, die ihm aus jener Aufgabe erwuchsen, bis ihm endlich — wie der Blitz einschlägt und ohne dass er sagen könne, wodurch — der befreiende Gedanke kam (30. August 1805, Nr. 123 T.). Die Lösung beruht einerseits auf der Umformung der Summe \(W_q\) in ein Produkt, andererseits auf den Eigenschaften der beiden eigentümlich gebauten Reihen
\[
f(x,m)=\sum_{i=0}^{\infty}(-1)^i\mathbin{.}(m,i),\quad F(x,m)=\sum_{i=0}^{\infty}x^{\frac{i}{2}}\mathbin{.}(m,i),
\]
wo
\[
(m,i)=\frac{(1-x^m)(1-x^{m-1})\ldots(1-x^{m-+1})}{(1-x)(1-x^2)\ldots(1-x^i)},
\]
und deren jede geeignet ist, jene Umformung zu leisten. Aus ihren Eigenschaften entspringt u. a. die für die additive Zahlentheorie bedeutsame Formel
\[
\prod_{i=1}^{\infty}\frac{1-x^{2i}}{1-x^{2i-1}}=1+\sum_{i=1}^{\infty}x^{\frac{i(i-1)}{2}},
\]
sowie für den Fall eines ungeraden \(n\) die Gleichung
\[
W_q=(r^q-r^{-q})(r^{2q}-r^{-2q})\ldots(r^{(n-2)q}-r^{-(n-2)q})
\]
und hieraus für ein ungerades \(n\)
\[
W_1=+\sqrt{(-1)^{\frac{n-1}{2}}\mathbin{.}n},
\]
während \(W_1=+(1+i)\sqrt{n}\) oder gleich Null ist, jenachdem \(n\equiv0\) oder \(n\equiv2\) (mod. 4). Ferner folgt, wenn \(n\) eine ungerade Primzahl ist, allgemein
\[
W_q=\left(\frac{q}{n}\right)\mathbin{.}\sqrt{(-1)^{\frac{n-1}{2}}\mathbin{.}n}.
\]
\newpage\SourceLeaf{54}\ArithmeticHead{47}{R}{DIE SPÄTEREN BEWEISE DES FUNDAMENTALTHEOREMS.}
\noindent Ist aber \(n=abc\ldots\) ein Produkt verschiedener ungerader Primzahlen, so besteht die Beziehung
\[
W_1=\prod_{a,b,c,\ldots}\biggl(1+r^{\frac{n^2}{a^2}}+r^{4\frac{n^2}{a^2}}+\ldots+r^{(a-1)^2\frac{n^2}{a^2}}\biggr),
\]
aus welcher folgender Satz hervorgeht: Die Anzahl der \(a,b,c,\ldots\), von denen bezw. \(\frac{n}{a},\frac{n}{b},\frac{n}{c},\ldots\) Nichtreste sind, ist gerade oder ungerade, je nachdem die Anzahl \(\nu\) der Zahlen \(a,b,c,\ldots\) welche die Form \(4\mu+3\) haben, kongruent \(0,1\) oder \(2,3\) (mod. 4) ist. Wird \(n\) als Produkt von nur zwei solchen Primzahlen gedacht, so folgt hieraus das Reziprozitätsgesetz. Ähnliche Betrachtungen ergeben auch die beiden Ergänzungssätze.

Im Anfang seiner Abhandlung bezeichnet \textsc{Gauss} die Summen \(W_q\) als eine reiche Quelle für Untersuchungen, deren Darstellung er für eine andere Stelle verheisst. Ob \textsc{Gauss} hier Beziehungen dieser Summen zu den \(\theta\)-Reihen gemeint habe, die nach aufgefundenen Stellen des Nachlasses\footnote{Siehe insbesondere W. III, S. 433 ff. und die Nachträge zur Analysis in W. X 1, S. 145 ff.; vergl. auch den weiter unten folgenden Aufsatz »Über Gauss’ Arbeiten zur Funktionentheorie« von L. \textsc{Schlesinger}.}\ ihm nicht unbekannt geblieben sind, stehe dahin. Wahrscheinlicher ist hier die mit art. 19 beginnende, aber schnell abbrechende Fortsetzung der \textit{Disquisitionum circa aequationes puras ulterior evolutio} (W. II, S. 263—65) gemeint, möglicherweise aber auch der zweite der beiden Beweise des Fundamentaltheorems, welche die dritte der oben genannten Abhandlungen enthält. Auch er entstammt der Kreisteilung, ist aber von \textsc{Gauss} in einer Form dargestellt, die ihn der Theorie der höheren Kongruenzen näher bringt. Bedeutet \(g\) eine primitive Wurzel für die ungerade Primzahl \(p\) und \(\xi\) den Ausdruck
\[
x-x^g+x^{g^2}-\ldots-x^{g^{p-2}},
\]
so lässt sich zeigen, dass \(\xi^2-(-1)^{\frac{p-1}{2}}\mathbin{.}p\), also auch, wenn \(q\) eine andere ungerade Primzahl bedeutet,
\[
\xi^{q-1}-(-1)^{\frac{p-1}{2}\mathbin{.}\frac{q-1}{2}}\mathbin{.}p^{\frac{q-1}{2}}
\]
durch \(\frac{1-x^p}{1-x}\) teilbar ist, desgleichen
\newpage\SourceLeaf{55}\ArithmeticHead{48}{L}{BACHMANN, ÜBER GAUSS’ ZAHLENTHEORETISCHE ARBEITEN.}
\[
x^q-x^{qg}+x^{qg^2}-\ldots-x^{qg^{p-2}}-\left(\frac{q}{p}\right)\mathbin{.}\xi.
\]
Durch Verbindung dieser Umstände mit der Kongruenz
\[
x^q-x^{qg}+x^{qg^2}-\ldots-x^{qg^{p-2}}\equiv\xi^q\ (\mathrm{mod.}\ q)
\]
wird man einfach zum Reziprozitätsgesetze geführt.
